\documentclass[11pt]{article}

\usepackage[a4paper,margin=1in]{geometry}
\usepackage{amsmath,amssymb,amsthm,bm}
\usepackage{booktabs}
\usepackage{array}
\usepackage{tabularx}
\usepackage{enumitem}
\usepackage{xcolor}
\usepackage{algorithm}
\usepackage{algpseudocode}
\usepackage{hyperref}
\usepackage{microtype}
\usepackage{xspace}

\hypersetup{
    colorlinks=true,
    linkcolor=blue!50!black,
    citecolor=blue!50!black,
    urlcolor=blue!50!black
}

\newtheorem{assumption}{Assumption}
\newtheorem{remark}{Remark}
\newtheorem{theorem}{Theorem}

\newcommand{\F}{\mathcal F}
\newcommand{\A}{\mathcal A}
\newcommand{\Lcal}{\mathcal L}
\newcommand{\R}{\mathbb R}
\newcommand{\E}{\mathbb E}
\newcommand{\Prob}{\mathbb P}
\newcommand{\norm}[1]{\left\lVert #1 \right\rVert}
\newcommand{\argmax}{\operatorname*{arg\,max}}
\newcommand{\argmin}{\operatorname*{arg\,min}}
\newcommand{\ACQ}{\textsc{Acquire}\xspace}
\newcommand{\TRACK}{\textsc{Track}\xspace}

% pseudocode readability
\algrenewcommand\algorithmicrequire{\textbf{Input:}}
\algnewcommand\algorithmicand{\textbf{and}}

\title{Recurrence-Aware Thompson Sampling\\
\large A Simple Algorithm for Bandits with Recurring Regimes}
\author{}
\date{}

\begin{document}
\maketitle

\begin{abstract}
We study a decision problem in which, at every round, a learner picks one of $K$ actions and
receives a reward. The catch is that the world is not fixed: it switches, from time to time,
between a small number of hidden ``regimes,'' and the \emph{same} regimes keep coming back.
A good learner should (i) quickly figure out which regime is currently active, (ii) exploit it,
(iii) notice when it changes, and (iv) \emph{remember} regimes it has already learned so that a
returning regime is recognized instead of relearned from scratch. We present a deliberately
minimal algorithm, \emph{Recurrence-Aware Thompson Sampling} (RA-TS), that does all four using
only three familiar pieces: Thompson sampling to choose actions, a change detector to spot when
a regime ends, and a small memory of past regimes to warm-start on their return. This note
explains the problem, the algorithm, and its guarantees in self-contained terms, states the
algorithm in one pseudocode block, and lists every parameter in one table.
\end{abstract}

\section{The problem, in plain terms}
\label{sec:problem}

There are $K$ \emph{actions} (also called arms), $\A=\{1,\dots,K\}$. Time proceeds in rounds
$t=1,2,\dots,T$. In round $t$ the learner chooses an action $A_t\in\A$ and observes a numerical
reward $Y_t$. The learner wants the total reward over $T$ rounds to be as large as possible.

The environment is \emph{piecewise stationary}: the timeline is cut into \emph{episodes}, and
inside one episode the rewards are governed by a single hidden \emph{regime} $z$. When an episode
ends, the regime changes. The key feature of our setting is \emph{recurrence}: only a few
distinct regimes ever occur, and they reappear over and over.

\paragraph{Level versus shape.}
Within a regime $z$, action $k$ has an (unknown) average reward
\begin{equation}
    \mu_k^{(z)} \;=\; \underbrace{\ell^{(z)}}_{\text{level}} \;+\; \underbrace{\theta_k^{(z)}}_{\text{shape}},
    \qquad \sum_{k=1}^K \theta_k^{(z)} = 0 .
    \label{eq:level-shape}
\end{equation}
The \emph{level} $\ell^{(z)}$ is a number added to every action equally (an overall ``how good
are things right now'' offset). The \emph{shape} $\bm\theta^{(z)}$ says which actions are better
or worse than average; it sums to zero by definition. The best action of a regime,
$k^\star(z)=\argmax_k \theta_k^{(z)}$, depends only on the shape. This distinction matters:
if every reward suddenly rises or falls by the same amount, the level changes but the best
action does not, so that is \emph{not} a real regime change for our purposes. A real regime
change is a change of \emph{shape}.

To compare shapes we use \emph{centering}. For any vector $\bm v\in\R^K$, subtract its own
average:
\begin{equation}
    \phi(\bm v) \;=\; \bm v - \frac1K\sum_{j=1}^K v_j .
\end{equation}
Centering removes the level automatically: $\phi(\bm\mu^{(z)})=\bm\theta^{(z)}$. Equivalently,
we can compare actions by their \emph{differences} $\mu_i-\mu_j$, which also do not depend on the
level. This is why the algorithm never needs to estimate $\ell$.

\paragraph{Noise.}
Observed rewards are the true mean plus zero-mean noise, $Y_t = \mu_{A_t}^{(z)} + \varepsilon_t$,
where $\varepsilon_t$ is $\sigma$-sub-Gaussian (informally: noise of typical size $\sigma$, with
light tails). This is the only probabilistic assumption needed to state the algorithm.

\paragraph{Goal.}
We measure performance by \emph{regret}: the total reward we lose compared with an oracle that
always plays the best action of the current regime. Section~\ref{sec:regret} bounds it.

\section{The idea in one paragraph}
\label{sec:idea}

RA-TS combines three standard tools. \textbf{(1)~Thompson sampling} chooses actions during normal
operation: it keeps a probability estimate of each action's mean and plays the one that looks
best under a random draw, which automatically balances trying and exploiting and needs no tuning.
\textbf{(2)~A change detector} watches the reward of the action we are playing and raises an alarm
when it drifts away from what the current regime predicts. \textbf{(3)~A memory (``library'')} of
regimes we have already learned lets us \emph{recognize} a returning regime from its shape and
\emph{resume} its learning instead of starting over. The only extra ingredient is a very short,
bounded ``probe'' that tries each action a few times right after a change, so that we have enough
information to recognize the regime. Everything else is plain Thompson sampling.

\section{What the algorithm remembers}
\label{sec:state}

\paragraph{The library.}
The library $\Lcal=\{1,\dots,J\}$ holds one entry per regime discovered so far. Entry $g$ stores,
for each action $k$, how many times that action has been played in regime $g$ and the sum of the
rewards seen:
\begin{equation}
    N_{g,k}\ (\text{count}), \qquad S_{g,k}\ (\text{reward sum}), \qquad
    \widehat\mu_{g,k}=\frac{S_{g,k}}{N_{g,k}}, \qquad
    \widehat{\bm\theta}_g=\phi(\widehat{\bm\mu}_g).
    \label{eq:lib}
\end{equation}
The counts and sums $(N_g,S_g)$ are the \emph{only} things stored. The mean
$\widehat{\bm\mu}_g$ and the shape $\widehat{\bm\theta}_g$ are \emph{derived on demand} from them,
never maintained separately: whenever they are read they simply reflect the current counts.

\paragraph{When $(N_g,S_g)$---and hence $\widehat{\bm\theta}_g$---change.}
This happens at exactly two kinds of event, and nowhere else:
\begin{enumerate}[leftmargin=1.8em,label=(\roman*),itemsep=1pt]
\item \emph{Once, when a probe is assigned to regime $g$}: the whole identification buffer is added
in one shot, $N_g \mathrel{+}= \bm n,\ S_g \mathrel{+}= \bm s$.
\item \emph{Every \TRACK round in regime $g$}: only the single action $a$ just played is
incremented, $N_{g,a}\mathrel{+}=1,\ S_{g,a}\mathrel{+}=Y$.
\end{enumerate}
So the shape estimate $\widehat{\bm\theta}_g$ is refreshed on \emph{every} \TRACK round (through
$\widehat{\bm\mu}_g$), yet it is \emph{read} in only one place: during \ACQ, when a freshly collected
probe buffer is matched against the library. The mean vector $\widehat{\bm\mu}_g$ is read once more,
at the single moment a regime is identified and \TRACK begins, to initialize the drift detector's
baseline $\mathrm{base}\gets\widehat{\bm\mu}_g$; that value is then frozen for the entire \TRACK
phase and does not follow the subsequent updates to $(N_g,S_g)$. Two consequences are worth noting. First, because
Thompson sampling plays mostly the best action, \TRACK sharpens mainly that one coordinate of
$\widehat{\bm\theta}_g$; the other coordinates keep the coarser accuracy they had from the probe.
Second, this causes no bias: each coordinate $\widehat\mu_{g,k}=S_{g,k}/N_{g,k}$ is an unbiased
estimate of action $k$'s mean no matter how few or how many times $k$ was played, so
$\widehat{\bm\theta}_g$ remains an unbiased shape estimate despite the very unequal counts. Its
usefulness for matching is therefore limited by the \emph{least}-sampled actions, i.e.\ by the
radius $b(n_{\min})$---exactly the quantity that appears in the analysis. These raw statistics are
also precisely the sufficient statistics Thompson sampling needs, so a regime is resumed just by
continuing to add to them.

\paragraph{The scratch buffer and the baseline.}
Two more short-lived quantities are kept for the current situation: a \emph{buffer} $(\bm n,\bm s)$
of counts and sums collected while identifying the regime, and, once a regime is being tracked, a
frozen \emph{baseline} vector $\mathrm{base}=\widehat{\bm\mu}_g$ (the regime's mean reward captured
at the moment the regime is identified and \TRACK begins) together with a per-action detector
statistic $G_k$.

\section{The two modes}
\label{sec:modes}

\paragraph{\ACQ (recognize the regime).}
Right after a change, we do not yet know which regime we are in. We run a short \emph{coverage
probe}: repeatedly play the least-tried action until every action has been tried $n_{\min}$ times.
This guarantees we have a few samples of \emph{every} action, which is exactly what is needed to
estimate the shape. The moment coverage is reached, we compute the buffer's shape
$\widehat{\bm\psi}=\phi(\bm s/\bm n)$ and compare it to every library entry using the distance
\begin{equation}
    D_g \;=\; \norm{\widehat{\bm\psi} - \widehat{\bm\theta}_g}_\infty
    \qquad (\text{largest per-action shape difference}).
\end{equation}
If the closest entry $g^\star$ is within a tolerance $\tau_{\mathrm{match}}$, we decide it is that
regime and \emph{warm-start}: we fold the buffer into entry $g^\star$ and continue its learning.
Otherwise the shape matches nothing we know, so we create a new library entry. Either way we then
switch to \TRACK. Because comparison uses centered shapes (equivalently, action differences), it
ignores the reward level entirely.

\paragraph{\TRACK (exploit and watch).}
Now that the regime is known, we run Thompson sampling on its stored posterior. For each action
we draw a random plausible mean and play the best draw:
\begin{equation}
    \widetilde\mu_k \sim \mathcal N\!\big(m_{g,k},\, v_{g,k}\big),
    \qquad A_t=\argmax_k \widetilde\mu_k,
    \label{eq:ts}
\end{equation}
where $(m_{g,k},v_{g,k})$ is the usual Gaussian posterior mean and variance built from the prior
$\mathcal N(\mu_0,\tau_0^2)$ and the stored data $(N_{g,k},S_{g,k})$:
\begin{equation}
    v_{g,k}=\Big(\tfrac{1}{\tau_0^2}+\tfrac{N_{g,k}}{\sigma^2}\Big)^{-1},
    \qquad
    m_{g,k}=v_{g,k}\Big(\tfrac{\mu_0}{\tau_0^2}+\tfrac{S_{g,k}}{\sigma^2}\Big).
    \label{eq:post}
\end{equation}
As more data accumulate, the draws concentrate and Thompson sampling plays the best action almost
always: deterministic exploitation appears on its own, with no separate ``commit'' step.

At the same time, we watch for drift. After playing action $a$ and seeing reward $Y_t$, we measure
how surprising it is relative to the frozen baseline and accumulate the surprise in a one-sided
CUSUM statistic:
\begin{equation}
    e_t = Y_t - \mathrm{base}[a],
    \qquad
    G_a \leftarrow \max\{0,\; G_a + |e_t| - \nu\}.
    \label{eq:cusum}
\end{equation}
An alarm is raised when $G_a \ge h$, and we return to \ACQ. Three deliberate choices make this
work: the baseline is \emph{frozen} (a running average would silently follow the drift and hide
it); the surprise uses the \emph{absolute} residual $|e_t|$ (so both increases and decreases are
caught); and the slack $\nu$ is set \emph{above} the average noise, so with no drift the statistic
is pushed back to zero (few false alarms), while a real shift makes it climb.

\section{The algorithm and its parameters}
\label{sec:alg}

Algorithm~\ref{alg:rats} states RA-TS in full. Table~\ref{tab:params} lists every parameter, what
it means, and how to set it.

\begin{algorithm}[t]
\caption{Recurrence-Aware Thompson Sampling (RA-TS)}
\label{alg:rats}
\begin{algorithmic}[1]
\Require actions $K$, noise scale $\sigma$, prior $(\mu_0,\tau_0)$, coverage $n_{\min}$,
match tolerance $\tau_{\mathrm{match}}$, detector slack $\nu$, threshold $h$
\Statex \emph{Throughout: $\widehat{\bm\mu}_g=S_g/N_g$ and $\widehat{\bm\theta}_g=\phi(\widehat{\bm\mu}_g)$ are recomputed from the current $(N_g,S_g)$.}
\State library $\Lcal \gets \varnothing$;\quad mode $\gets \ACQ$;\quad buffer $(\bm n,\bm s)\gets(\bm 0,\bm 0)$
\For{$t = 1,2,\dots,T$}
  \If{mode $=\ACQ$}
     \State play $a \gets \argmin_k n_k$, observe $Y$ \Comment{coverage probe: try the least-sampled action}
     \State $n_a \gets n_a + 1$;\quad $s_a \gets s_a + Y$
     \If{$\min_k n_k \ge n_{\min}$} \Comment{every action covered: identify the regime}
        \State $\widehat{\bm\psi} \gets \phi(\bm s / \bm n)$;\quad
               $g^\star \gets \argmin_{g\in\Lcal} \norm{\widehat{\bm\psi}-\widehat{\bm\theta}_g}_\infty$
        \If{$\Lcal \ne \varnothing$ \algorithmicand{} $\norm{\widehat{\bm\psi}-\widehat{\bm\theta}_{g^\star}}_\infty \le \tau_{\mathrm{match}}$}
           \State $N_{g^\star} \gets N_{g^\star} + \bm n$;\quad $S_{g^\star} \gets S_{g^\star} + \bm s$;\quad $g \gets g^\star$ \Comment{match: warm-start}
        \Else
           \State append new entry with $(N,S) \gets (\bm n,\bm s)$;\quad $g \gets$ new index \Comment{new regime}
        \EndIf
        \State mode $\gets \TRACK$;\quad $\mathrm{base} \gets \widehat{\bm\mu}_g$;\quad $\bm G \gets \bm 0$
     \EndIf
  \Else \Comment{mode $=\TRACK$}
     \State draw $\widetilde\mu_k \sim \mathcal N(m_{g,k},v_{g,k})$ for all $k$ \Comment{Eq.~\eqref{eq:post}}
     \State play $a \gets \argmax_k \widetilde\mu_k$, observe $Y$ \Comment{Thompson sampling}
     \State $S_{g,a} \gets S_{g,a} + Y$;\quad $N_{g,a} \gets N_{g,a} + 1$ \Comment{resume/update posterior}
     \State $G_a \gets \max\{0,\; G_a + |Y - \mathrm{base}[a]| - \nu\}$ \Comment{drift CUSUM on played action}
     \If{$G_a \ge h$} \Comment{alarm: regime ended}
        \State mode $\gets \ACQ$;\quad buffer $(\bm n,\bm s) \gets (\bm 0,\bm 0)$
     \EndIf
  \EndIf
\EndFor
\end{algorithmic}
\end{algorithm}

\begin{table}[t]
\centering
\small
\renewcommand{\arraystretch}{1.35}
\begin{tabularx}{\textwidth}{@{}c >{\raggedright\arraybackslash}X >{\raggedright\arraybackslash}p{0.36\textwidth}@{}}
\toprule
\textbf{Symbol} & \textbf{Meaning} & \textbf{How to set (and typical value)} \\
\midrule
$K$ & number of actions & problem-given \\
$\sigma$ & noise scale (sub-Gaussian proxy) & estimated from data \\
$\delta$ & failure probability of the confidence bounds & small, e.g.\ $0.1$; enters $L=2\log(8KT/\delta)$ \\
$T$ & horizon (number of rounds) & problem-given; used only inside $L$ \\
$\mu_0$ & prior mean of the Thompson posterior & the average reward level; not sensitive \\
$\tau_0$ & prior standard deviation (diffuse) & large relative to reward spread, e.g.\ $0.5$ \\
$n_{\min}$ & samples per action required by the probe & $n_{\min} > 32\,\sigma^2 L/\Delta_{\mathrm{sh}}^2$; small in practice (e.g.\ $3$) \\
$\tau_{\mathrm{match}}$ & shape-match tolerance & in $[\,b(n_{\min}),\,\Delta_{\mathrm{sh}}-b(n_{\min})\,]$; e.g.\ $0.8\,\Delta_{\mathrm{sh}}$ \\
$\nu$ & detector slack (per-step ``ignore'' band) & $\sigma\sqrt{2/\pi} < \nu < \eta_{\min}$; e.g.\ $1.3\,\sigma$ \\
$h$ & detector alarm threshold & larger $=$ fewer false alarms, slower detection; e.g.\ $8\,\sigma$ \\
\midrule
\multicolumn{3}{@{}l}{\emph{Derived (not tuned):}\ \
$r(n)=\sigma\sqrt{L/n}$,\ \ $b(n)=2\,r(n)$,\ \
$\phi(\bm v)=\bm v-\tfrac1K\sum_j v_j$.} \\
\bottomrule
\end{tabularx}
\caption{All parameters of RA-TS. $\Delta_{\mathrm{sh}}$ is the smallest shape separation between
distinct regimes and $\eta_{\min}$ the smallest change of the played action's mean at a regime
switch (Assumptions~\ref{ass:sep}--\ref{ass:detect}); both are properties of the environment, not
tuned. The radius $r(n)$ is the size of the estimation error after $n$ samples; $b(n)$ is its
centered (shape) counterpart.}
\label{tab:params}
\end{table}

\section{Why each piece is there}
\label{sec:why}

\paragraph{Why a probe, and not just Thompson sampling?}
It is tempting to drop the probe and let Thompson sampling itself explore during \ACQ. This does
not work, for a simple reason. Recognizing a regime means telling the actions apart, which
requires at least a few samples of \emph{every} action, including the bad ones. But Thompson
sampling is built to \emph{stop} playing actions that look bad: after one or two samples a clearly
inferior action is almost never chosen again. So coverage is never reached, identification never
happens, the algorithm stays in \ACQ, and---since the detector only runs in \TRACK---a later
regime change is not even watched. The probe supplies exactly the missing samples, and nothing
more: at most $n_{\min}K$ plays, once per episode.

\paragraph{Why remember regimes?}
Without the library, every reappearance of a regime would be learned from scratch, paying the full
exploration cost each time. With the library, a returning regime is recognized from its shape and
its posterior is resumed, so its learning cost is paid essentially once and then amortized over all
future visits. This is the main advantage under frequent switching among a fixed set of regimes.

\paragraph{Why watch only the played action?}
The detector updates $G_a$ only for the action actually played, which is almost always the current
best action. This is enough to catch any change that moves the best action's reward, and it costs
nothing extra. The price is a blind spot: a change that alters only actions we are \emph{not}
playing cannot be seen. This is unavoidable for any method that does not periodically re-test every
action.

\section{Guarantees}
\label{sec:regret}

We now state what can be proved. We first list the assumptions in words, then give the regret
bound and a sketch of why each part holds.

\begin{assumption}[Separated regimes]
\label{ass:sep}
Distinct regimes differ in shape by at least $\Delta_{\mathrm{sh}}$: for any two regimes that
occur, $\norm{\bm\theta^{(z)}-\bm\theta^{(z')}}_\infty\ge\Delta_{\mathrm{sh}}>0$. (Regimes are
actually distinguishable.)
\end{assumption}

\begin{assumption}[Visible changes]
\label{ass:visibility}
When a regime ends, the mean of the action we are committed to moves by at least $\eta_{\min}$.
(A change shows up on the action we are actually playing.)
\end{assumption}

\begin{assumption}[Feasible detector and probe]
\label{ass:detect}
The constants satisfy $\sigma\sqrt{2/\pi} < \nu < \eta_{\min}$ (the slack sits above the noise but
below the visible change) and $n_{\min} > 32\,\sigma^2 L/\Delta_{\mathrm{sh}}^2$, i.e.\
$b(n_{\min}) < \Delta_{\mathrm{sh}}/4$ (a probe resolves the shape to finer than the regime gap),
with $\tau_{\mathrm{match}}\in[\,b(n_{\min}),\,\Delta_{\mathrm{sh}}-b(n_{\min})\,]$.
\end{assumption}

\begin{theorem}[Regret bound, high probability]
\label{thm:main}
Let $G_T$ be the number of episodes, $\mathcal R$ the set of distinct regimes, $T_r$ the total
number of \TRACK rounds spent in regime $r$ (so $\sum_r T_r\le T$), $\Delta_{r,k}$ the reward gap
of action $k$ in regime $r$, and $S_T$ the number of ``short'' episodes (shorter than the
detection delay plus the probe). Under Assumptions~\ref{ass:sep}--\ref{ass:detect}, with
probability at least $1-\delta$,
\begin{equation}
   R_T \;\le\;
   \underbrace{\sum_{r\in\mathcal R}\ \sum_{k\ne k^\star(r)} \frac{C\,\sigma^2 \log T_r}{\Delta_{r,k}}}_{\text{(i)}}
   \;+\;
   \underbrace{\Delta_{\max}\,(G_T + S_T)\Big(n_{\min}K + \tfrac{h}{\eta_{\min}-\nu}\Big)}_{\text{(ii)}}
   \;+\;
   \underbrace{\delta\,T\,\Delta_{\max}}_{\text{(iii)}},
   \notag
\end{equation}
for a universal constant $C$, where $\Delta_{\max}$ bounds the one-round regret.
\end{theorem}

Each term corresponds to one thing the algorithm does:
\begin{itemize}[leftmargin=1.6em,itemsep=1pt]
\item \textbf{(i) Learning.} The cost of Thompson sampling finding the best action in each
regime---paid \emph{once per regime} and shared over all its recurrences (note the $\log$ of the
cumulative lifetime $T_r$).
\item \textbf{(ii) Identify {}+{} detect.} A fixed overhead per episode: the bounded probe
($\le n_{\min}K$ plays) plus the detection delay ($\approx h/(\eta_{\min}-\nu)$ rounds). The
$S_T$ part is the small extra cost on ``short'' episodes---those shorter than this overhead, where
a change can land inside the probe.
\item \textbf{(iii) Failure.} The rare event that a confidence bound is violated.
\end{itemize}

\paragraph{Sketch of the learning term.}
All confidence estimates hold simultaneously on a ``good event'' of probability $\ge 1-\delta$
(a union bound over actions and sample sizes, which is why $\log(8KT/\delta)$ appears in the
radii). On that event, because a regime's posterior is only ever \emph{added to} and never reset,
all Thompson plays ever made in regime $r$ behave like one uninterrupted Thompson run of length
$T_r$ on a fixed problem. The standard finite-time Thompson-sampling bound for such a run is
$\sum_{k\ne k^\star}O(\sigma^2\log T_r/\Delta_{r,k})$. The important point is that the logarithm is
of the \emph{cumulative} lifetime $T_r$, so revisiting a regime does not restart the exploration
cost---this is the recurrence benefit.

\paragraph{Sketch of the probe and delay term.}
Each episode triggers at most one \ACQ: a probe of at most $n_{\min}K$ forced plays, plus the
rounds between the true change and the alarm. After a change the residual $|e_t|$ is at least
$\eta_{\min}$ on average, so the CUSUM climbs at rate $\ge\eta_{\min}-\nu$ and crosses $h$ in about
$h/(\eta_{\min}-\nu)$ rounds. Each such round costs at most $\Delta_{\max}$.

\paragraph{Sketch of correct identification and the short-episode term.}
After a clean probe, the shape estimate is within $b(n_{\min})$ of the truth
(Assumption~\ref{ass:detect}), which is closer than half the regime gap, so the nearest library
entry is the right regime (if known) and otherwise a new entry is created---identification is
correct. It can fail only when a change lands \emph{inside} the short probe window, mixing two
regimes; this happens only on episodes shorter than the probe-plus-delay overhead, and each such
case costs at most one extra \ACQ cycle. Finally, choosing $\nu>\sigma\sqrt{2/\pi}$ makes the
detector's ``no drift'' increments negative on average, so false alarms are exponentially rare and
contribute only to the $\delta T$ failure budget.

\paragraph{Reading the bound.}
Only the first term depends on the gaps, and it is paid once per regime and then shared across all
recurrences. The remaining terms are a small, fixed overhead per episode (identify $+$ detect),
with a little extra only for episodes too short to finish identifying. In the regime of
interest---a few regimes recurring many times, each episode comfortably longer than the detection
delay---the bound is dominated by the amortized Thompson term.

\end{document}
